\documentclass[11pt,a4paper]{article}
\usepackage[margin=1in]{geometry}
\usepackage{hyperref}
\usepackage{enumitem}
\usepackage{graphicx}
\usepackage{xcolor}
\usepackage{titling}

% -----------------------------------------------------
% UCS503P Project Proposal formatting
% -----------------------------------------------------

% Variable: Lab Instructor
\makeatletter
\newcommand{\BvrSetLabInstructor}[1]{%
  \gdef\Bvr@LabInstructor{#1}}
\newcommand{\Bvr@LabInstructor}{Ms. Nisha Thakur}
\newcommand{\BvrLabInstructorValue}{\Bvr@LabInstructor}
\makeatother

% Title and subtitle formatting
\pretitle{\begin{center}\bfseries\fontsize{18bp}{18bp}\selectfont}
\makeatletter
\newcommand{\BvrSubtitle}[1]{%
  \posttitle{%
    \par\end{center}
    \begin{center}\large#1\end{center}
    \vskip 0.5em}%
}
\makeatother

% Hints in section titles
\newcommand{\BvrHint}[1]{%
  {\normalfont\normalsize\itshape\hfill #1}}

% -----------------------------------------------------

\title{VR Everywhere - Virtual Travel Experience}

\BvrSetLabInstructor{Ms. Nisha Thakur}

\BvrSubtitle{%
  Project Proposal for UCS503P \\
  Submitted to: Ms. Nisha Thakur \\
  \bfseries Thapar Institute of Engineering and Technology%
}

\author{
  Abishek Garg, CSE
  {Roll No: \texttt{1024170428}}
  \and
  Deepak, CSE
  {Roll No: \texttt{1024170432}}
  \and
  Harnoor Kaur Dran, CSE
  {Roll No: \texttt{1024170436}}
}

\date{July -- December 2026}

\begin{document}
\maketitle

\section{Higher-order goal%
  \BvrHint{\ldots immersive travel experience}}

This project aims to make travel exploration more accessible by allowing users to virtually
experience tourist destinations without physically travelling to them. While the broader goal is to
explore the potential of virtual reality in tourism and education, the immediate engineering
objective is to develop a simple, deployable VR-based travel application that allows users to explore
selected destinations through immersive environments.

The project will focus on providing an easy-to-use interface, destination information, interactive
exploration, and basic VR navigation rather than attempting to create a complete commercial
tourism platform.

\section{Time-to-value%
  \BvrHint{\ldots primary heuristic}}

\begin{itemize}[leftmargin=0.7cm]
    \item We will deliver meaningful increments quickly by first implementing the core workflow:
    selecting a destination $\rightarrow$ viewing its information $\rightarrow$ starting a VR tour
    $\rightarrow$ exploring the environment.
    \item We will prioritize a small number of high-quality destinations over attempting to build a large virtual world.
    \item Development will follow short weekly iterations, with each iteration producing a demonstrable feature or improvement.
    \item Testing will be performed continuously during development rather than being postponed until the final stage.
    \item Success will initially be defined by whether a user can successfully discover a destination
    and complete a basic VR tour, followed by improvements based on user feedback.
\end{itemize}

\section{Problem Statement}

Travelling to tourist destinations is not always possible for everyone due to cost, time constraints,
physical limitations, geographical distance, or other circumstances.

\begin{itemize}[leftmargin=0.7cm]
    \item \textbf{Accessibility:} People who cannot physically travel may have limited opportunities to experience tourist destinations.
    \item \textbf{Cost:} International and long-distance travel can be expensive.
    \item \textbf{Time constraints:} Students and working individuals may not have sufficient time to visit multiple destinations.
    \item \textbf{Limited pre-travel experience:} Traditional travel websites provide photographs and videos, but they do not provide an immersive experience of the location.
    \item \textbf{Information fragmentation:} Information about a destination, its history, attractions, and visual experience is often spread across multiple platforms.
\end{itemize}

Virtual Reality provides an opportunity to address these limitations by allowing users to experience
selected tourist destinations in an immersive environment.

\section{Proposed Solution%
  \BvrHint{\ldots Software Proposition}}

\subsection{Overview}

Build a VR-based ``Virtual Travel Experience'' system with the following components:

\begin{itemize}[leftmargin=0.7cm]
    \item \textbf{Destination explorer:} Users can browse and search for available tourist destinations.
    \item \textbf{Destination information:} Users can view basic information such as location, description, historical facts, and important attractions.
    \item \textbf{VR tour:} Users can enter a virtual environment and explore a selected destination.
    \item \textbf{Interactive hotspots:} Important locations within the VR environment can be selected to display additional information.
    \item \textbf{Audio guide:} Basic audio narration can be provided for selected attractions.
    \item \textbf{Favorites:} Users can save destinations for later exploration.
    \item \textbf{User profile:} Users can view their saved and previously explored destinations.
\end{itemize}

The initial implementation will contain a limited number of destinations to keep the project
achievable within the three-month development period.

\subsection{Core workflow}

\begin{enumerate}[leftmargin=0.7cm]
    \item User opens the application and logs in or continues as a user.
    \item User views available destinations on the Home/Explore screen.
    \item User selects a destination.
    \item The system displays destination information and provides a Start VR Tour option.
    \item User enters the VR environment and explores the destination.
    \item User can interact with selected hotspots to view information or listen to an audio guide.
    \item User exits the tour and can save the destination to Favorites.
\end{enumerate}

This keeps the central workflow simple and demonstrable.

\subsection{Operational constraints for an educational project}

\begin{itemize}[leftmargin=0.7cm]
    \item The application will initially support a small number of destinations, approximately 3--5, rather than attempting to model the entire world.
    \item VR environments will use available 360$^\circ$ images/assets or lightweight 3D environments instead of developing highly detailed environments from scratch.
    \item The application should remain usable on supported VR hardware and provide a basic non-VR/desktop mode where feasible for demonstration.
    \item Advanced features such as multiplayer VR, live tour guides, complex AI systems, and real-time environmental simulation are outside the initial scope.
    \item The system will use commonly available development tools and cloud hosting services.
\end{itemize}

\section{Solution Approach%
  \BvrHint{\ldots Engineering focus}}

This is an engineering course project, so the primary focus will be on system design, VR interaction,
usability, integration, testing, and deployment, rather than developing new VR algorithms or
conducting research into virtual reality.

\subsection{VR environment and interaction%
  \BvrHint{\ldots practical implementation}}

\begin{itemize}[leftmargin=0.7cm]
    \item Basic VR navigation.
    \item 360$^\circ$ environment viewing.
    \item Interactive information hotspots.
    \item Audio narration.
    \item Simple VR menu.
    \item Destination exit/return functionality.
\end{itemize}

The project will prioritize usability and stable interaction over highly complex graphical effects.

\subsection{Application architecture%
  \BvrHint{\ldots web-connected VR solution}}

A typical architecture will consist of:

\begin{itemize}[leftmargin=0.7cm]
    \item \textbf{VR Client:} Unity-based application responsible for the VR environment, navigation, interaction, and user interface.
    \item \textbf{Backend API:} Responsible for authentication, destination information, user preferences, favorites, and tour history.
    \item \textbf{Database:} Stores users, destinations, hotspot information, favorites, and tour history.
    \item \textbf{Cloud/Deployment Layer:} Hosts the backend and database and provides access to the application services.
\end{itemize}

The architecture separates the VR application from backend services so that destination information
can be updated without modifying the complete application.

\subsection{Development and deployment%
  \BvrHint{\ldots fast delivery enabler}}

\begin{itemize}[leftmargin=0.7cm]
    \item Use Git and GitHub for version control.
    \item Automatically build and test the backend when changes are pushed, where CI/CD is introduced.
    \item Run backend unit tests.
    \item Maintain a consistent deployment process.
    \item Package the VR application as a suitable build for the selected demonstration platform.
\end{itemize}

\section{Evaluation Criterion%
  \BvrHint{\ldots measurable and attributable}}

\subsection{Primary evaluation metric}

\textbf{Tour Completion Rate:} Percentage of test users who can successfully select a destination,
enter the VR tour, interact with a hotspot, and exit the tour without assistance.

\begin{itemize}[leftmargin=0.7cm]
    \item Target: At least 80\% of test users should successfully complete the workflow during user testing.
\end{itemize}

\subsection{Secondary evaluation metrics}

\begin{itemize}[leftmargin=0.7cm]
    \item \textbf{Usability:} Average user rating of the application interface on a 1--5 scale.
    \item \textbf{Navigation success:} Percentage of users who can successfully navigate or interact within the VR environment.
    \item \textbf{Performance:} Maintain a stable and comfortable VR experience on the selected hardware.
    \item \textbf{Feature usage:} Number of users who interact with information hotspots or audio guides.
    \item \textbf{Reliability:} Successful loading of destinations and VR environments during testing.
\end{itemize}

\subsection{Pilot validation plan%
  \BvrHint{\ldots high-level}}

\begin{itemize}[leftmargin=0.7cm]
    \item Recruit approximately 10--20 student users for testing.
    \item Provide users with a simple task such as: ``Find the Taj Mahal, start the VR tour, explore the location, open an information hotspot, and exit the tour.''
    \item Record whether users complete the task successfully.
    \item Collect feedback using a short questionnaire.
    \item Identify usability problems and improve the interface before final deployment.
\end{itemize}

\section{Scalability%
  \BvrHint{\ldots with foresight}}

\subsection{Scalable theoretically}

\begin{itemize}[leftmargin=0.7cm]
    \item Store destination information in the database.
    \item Use separate VR scenes/environments for destinations.
    \item Separate the frontend/VR client and backend services.
    \item Use indexed database queries for destination searches.
    \item Keep backend APIs stateless where practical.
\end{itemize}

\subsection{Operationally deployable}

\begin{itemize}[leftmargin=0.7cm]
    \item Cloud-hosted backend.
    \item Managed relational database.
    \item Version-controlled source code.
    \item Repeatable deployment process.
\end{itemize}

This allows the initial 3--5 destinations to eventually grow into a larger virtual tourism catalogue.

\section{Engine availability heuristic}

\begin{itemize}[leftmargin=0.7cm]
    \item \textbf{VR/Game Engine:} Unity
    \item \textbf{Programming Language:} C\#
    \item \textbf{VR Framework:} OpenXR + XR Interaction Toolkit
    \item \textbf{Backend:} Spring Boot
    \item \textbf{Backend Language:} Java
    \item \textbf{Database:} PostgreSQL
    \item \textbf{Authentication:} Spring Security + JWT
    \item \textbf{API:} REST API
    \item \textbf{API Testing:} Postman
    \item \textbf{Version Control:} Git + GitHub
    \item \textbf{Deployment:} Render / Railway / similar cloud platform
\end{itemize}

These technologies allow the project to focus on VR interaction, application workflow, testing, and
deployment rather than developing infrastructure from scratch.

\section{Project scope and deliverables%
  \BvrHint{\ldots iteration-friendly}}

\subsection{Initial deliverable%
  \BvrHint{\ldots first iteration}}

\begin{itemize}[leftmargin=0.7cm]
    \item Basic application interface.
    \item Home/Explore screen.
    \item Destination listing.
    \item Destination details page.
    \item One working VR destination.
    \item Basic VR navigation.
    \item At least one interactive information hotspot.
    \item Basic backend API.
    \item Database containing destination information.
    \item Git repository and basic testing.
\end{itemize}

\subsection{Subsequent deliverables}

\begin{itemize}[leftmargin=0.7cm]
    \item Additional VR destinations.
    \item Search functionality.
    \item Audio guide.
    \item User profile.
    \item Tour history.
    \item Improved VR interface.
    \item Backend integration.
    \item System and user testing.
    \item Performance optimization.
\end{itemize}

\section{Risks and mitigations}

\begin{center}
\begin{tabular}{p{0.31\textwidth} p{0.58\textwidth}}
\textbf{Risk} & \textbf{Mitigation} \\ \hline
VR hardware availability &
Develop and test on the available VR device and maintain a desktop mode where feasible for demonstration. \\[4pt]

Complexity of VR development &
Limit the initial project to 3--5 destinations and basic interactions. \\[4pt]

Performance issues &
Use optimized 3D assets, appropriate texture sizes, and lightweight environments. \\[4pt]

Limited development time &
Follow weekly iterations and prioritize the core VR tour workflow. \\[4pt]

Asset availability &
Use appropriately licensed/free assets and 360$^\circ$ environments rather than attempting to create every asset from scratch. \\[4pt]

User discomfort &
Keep movement simple, provide comfortable navigation options, and test the application with users. \\[4pt]

Deployment difficulties &
Deploy backend services early rather than waiting until the final week.
\end{tabular}
\end{center}

\section{Summary}

This project proposes a deployable VR Travel Experience platform that allows users to explore
selected tourist destinations through immersive virtual environments. The system will combine a
simple destination discovery interface with VR tours, interactive hotspots, basic destination
information, and user favorites.

The project meets the course requirements by emphasizing rapid development, structured software
engineering, measurable user evaluation, testing, and deployment. The initial scope is intentionally
limited to a small number of destinations and essential VR interactions so that a functional and
polished system can be completed within three months.

The architecture will also allow additional destinations and features to be introduced in future
iterations without redesigning the complete system.

\section*{Technology Stack}

\begin{center}
\begin{tabular}{p{0.30\textwidth} p{0.55\textwidth}}
\textbf{Component} & \textbf{Technology} \\ \hline
VR Application & Unity \\
Programming & C\# \\
VR Framework & OpenXR + XR Interaction Toolkit \\
Backend & Spring Boot \\
Backend Language & Java \\
Database & PostgreSQL \\
Authentication & Spring Security + JWT \\
API & REST API \\
API Testing & Postman \\
Version Control & Git + GitHub \\
Deployment & Render / Railway / similar cloud platform
\end{tabular}
\end{center}

\section*{Core User Interface}

\begin{itemize}[leftmargin=0.7cm]
    \item The main application navigation will contain five simple options:
    \textbf{Home | Explore | VR Tours | Favorites | Profile}
    \item Inside the VR tour, a lightweight floating menu will provide:
    \textbf{Map | Information | Audio | Settings | Exit Tour}
\end{itemize}

\end{document}
